\section{Lc incidence and completion of the bounded-fibre argument}
\label{sec:incidence}

It remains to prove part~(3) of Theorem~\ref{thm:bounded-fibre}.
We first consider a divisorial valuation $w$ whose restriction $u$ to $K$
is nonzero and satisfies $\bar A_\sigma(u)=0$.  Make the bounded parameter
and Kummer extensions of Sections~\ref{sec:bounded-presentations} and
\ref{sec:traits}, and prolong $w$ with its scale unchanged.  The barred
equality and the hypothesis on the zero set of $A_\sigma$ are preserved
under these extensions.  We use the notation of Proposition~\ref{prop:trait-comparison}.
In particular $s$ is the geometric closed point on the parameter base,
$b$ is its vector of boundary orders, and a label is an irreducible
component of a stratum in the \emph{ordinary} fibre $Y_s$.

Define the nonnegative piecewise linear function
\begin{equation}\label{inc:gap}
 g(x)=\phi(x)-\mu(p_*x).
\end{equation}
The subdivisions already made ensure that $g$ is linear on every source
cone.  Equations~\eqref{trait:T} and \eqref{trait:TE} give
\begin{equation}\label{inc:excess}
 A_\sigma(w)=g(x)+E(w),\qquad p_*x=b,
 \qquad E(w)\geq 0.
\end{equation}
Here $x$ is the boundary-weight vector of $w$, and the support cone
$\delta$ of $x$ carries its label.  We will connect this labelled
vector to a zero of $g$ in the same slice, with path length controlled
by $g(x)$.  The next two subsections establish the incidence of
zero-gap strata needed for that construction.

\subsection{An effective dlt model for the parameter family}

The incidence argument takes place on the algebraic parameter family
$f:Y\to B$, including the local Kummer changes, after extension to the
algebraically closed residue field $\kappa$.  It does not take place on
the formal trait.  The base is smooth and may be shrunk to a connected
affine neighbourhood of $s$.  The morphism is projective, its generic
fibre is geometrically integral, and its fibres are connected by Stein
factorization and normality of $B$.  Its generic fibre maps birationally
to the original effective lc form model, which we denote by $G_{\rm gen}$.
These properties persist after a connected pointed \'{e}tale base change.

Let $\partial Y$ be the reduced toroidal boundary, and pull all the
parameter data to the present base.  Consider the ordinary subpair
\begin{equation}\label{inc:subpair}
 C^\natural=\partial Y-
 \frac{1}{r_0}\operatorname{div}^{\log}_{Y/B}(\eta)+f^*M_0.
\end{equation}
The relative log determinant gives
$K_Y+C^\natural\sim_{\Q,B}0$.  The ordinary log discrepancy of a
toroidal ray is $g$ on that ray, including after subdivision.  Thus
$(Y,C^\natural)$ is sub-lc: on a smooth toroidal resolution its boundary
is SNC with coefficients at most one.  On the generic fibre it is
crepant to the effective form boundary on $G_{\rm gen}$.

The construction uses Birkar's very-exceptional-divisor method
\cite[Section~3 and Theorem~3.4]{BirkarFlips}; the argument below
specifies the finite model on which its negativity step is applied.

\begin{lemma}\label{inc:effective-model}
There is a smooth projective toroidal resolution $P\to Y$ and a finite
sequence of ordinary $\Q$-factorial dlt MMP steps over $B$ from $P$ to a
model $P_j$ with an effective dlt boundary $D_j$ such that
\[
 K_{P_j}+D_j\sim_{\Q,B}0,
\]
and $(P_j,D_j)$ is crepant to $(Y,C^\natural)$.  The lc centres of the
crepant subpair on $P$ and of $(P_j,D_j)$ correspond by strict transform,
preserving their inclusions
and their proper images in $B$.
\end{lemma}

\begin{proof}
Choose $P$ smooth with strict SNC toroidal boundary, and let
$C_P^\natural$ be the crepant transform of $C^\natural$.  Set
\begin{equation}\label{inc:truncation}
 D=\max\{C_P^\natural,0\},\qquad
 N=D-C_P^\natural.
\end{equation}
Then $(P,D)$ is effective dlt, $N\geq0$, and
\begin{equation}\label{inc:initial-adjoint}
 K_P+D\sim_{\Q,B}N,
 \qquad A_{C^\natural}(v)=A_D(v)+v(N).
\end{equation}
The lc places of the two pairs agree.  To check the nontrivial direction,
the centres of lc places of the SNC pair $(P,D)$ are intersections of
coefficient-one components, and none is contained in the components
truncated in \eqref{inc:truncation}.  Every such place has $v(N)=0$.
Conversely $A_{C^\natural}=0$ forces both nonnegative terms on the
right of \eqref{inc:initial-adjoint} to vanish.  On the generic fibre
$N$ is exceptional over $G_{\rm gen}$, since the form boundary on that
model is effective.

We construct an MMP with scaling, give the scaling alternative, and
then show that the exceptional error vanishes on a finite model.  All
the varieties used here are of finite type over $\kappa$.  If needed
for the ordinary complex statements of the MMP theorems, transport
the entire family by an abstract isomorphism $\kappa\simeq\C$.
Such an isomorphism exists because $\kappa$ is the algebraic closure of
a finitely generated extension of $\C$ and has the same transcendence
degree over $\Q$ as $\C$.  This identification need not fix the original
copy of $\C$.

Choose a rational ample divisor $H$ with $K_P+D+H$ nef over $B$.
Write $P_i,D_i,N_i,H_i$ for successive models and pushforwards, and
put $C_i^\natural=D_i-N_i$.  These subpairs are crepant throughout:
the adjoint of $C^\natural$ is rationally linearly pulled back from
the base, so, with compatible canonical divisors, its pullbacks to a
common resolution agree.  For each negative MMP step, with common
resolution maps $a$ and $a^+$, negativity therefore gives
\begin{equation}\label{inc:decreasing-pulls}
 a^*N_i-(a^+)^*N_{i+1}\geq0.
\end{equation}
All $N_i$ remain effective.  In particular
$K_{P_i}+D_i\sim_{\Q,B}N_i$ is relatively pseudo-effective, so a Mori
fibre ending is impossible.

\smallskip\noindent\emph{Scaling construction.}
Here are the existence and scaling details.  For any rational
$0<t\leq1$, choose a sufficiently small rational $\epsilon>0$.
The divisor $tH+\epsilon D$ is ample, and a sufficiently divisible
general representative $T_{t,\epsilon}$ gives an effective big klt
boundary
\begin{equation}\label{inc:klt-representative}
 \Delta_t=(1-\epsilon)D+T_{t,\epsilon}
 \sim_{\Q,B}D+tH.
\end{equation}
The general representative has arbitrarily small coefficients and
meets the SNC support transversely; $(P,(1-\epsilon)D)$ is klt.
If the steps so far have scaling thresholds at least $t$, they are
nonpositive for $K_P+\Delta_t$.  Hence its pushed boundary is still
klt, big, and rationally linearly equivalent to $D_i+tH_i$.
Bigness is preserved under these non-extractive birational maps.
On any such model a klt big boundary can moreover be represented
with a small general ample part.  Indeed take an effective big
decomposition with the required ample part and mix it, with sufficiently
small positive coefficient, into the existing klt representative.

Suppose that $K_{P_i}+D_i$ is not nef, while
$K_{P_i}+D_i+\lambda_{i-1}H_i$ is nef.  Its nef threshold
$\lambda_i$ on this interval is positive.  Choose a rational
$t\in(0,\lambda_i)$.  The klt cone theorem applied to a representative
at $t$ with ample part leaves only finitely many negative extremal rays
to test.  The remaining part of the cone is nonnegative both at $t$
and at $\lambda_{i-1}$, hence throughout their interval.  Thus
$\lambda_i$ is a rational maximum of finitely many ray thresholds,
and is attained on a ray $R_i$ with
\[
 (K_{P_i}+D_i)\cdot R_i<0,
 \qquad (K_{P_i}+D_i+\lambda_iH_i)\cdot R_i=0.
\]
The klt contraction theorem at $t$ contracts $R_i$.  At $\lambda_i$
the nef divisor is relatively semiample: use its klt representative
with an ample part and apply basepoint-freeness to that representative
with the ample part removed.  It therefore descends through the
contraction.  If the contraction is small, the klt flip supplied by
\cite{BCHM} at $t$ is also the flip for $K_{P_i}+D_i$: over the
contraction their classes are positive multiples, because the
$\lambda_i$ combination descends.  These are the ordinary dlt MMP
steps; they preserve $\Q$-factoriality and dlt, and are nonpositive
for all the positive-scaling combinations with parameter at most
$\lambda_i$.  This constructs the sequence until nefness, or constructs
an infinite sequence with nonincreasing thresholds.

\smallskip\noindent\emph{Limiting threshold.}
We next prove that an infinite sequence must have
\begin{equation}\label{inc:scaling-zero}
 \lambda_i\longrightarrow0.
\end{equation}
Every divisorial contraction lowers the relative Picard number, so
there are only finitely many.  Fix an index $j$ after the last one.
If the thresholds had positive limit, choose a rational closed
interval bounded away from zero, with upper endpoint at most the previous
scaling bound, containing all sufficiently late thresholds.  On $P_j$, choose
klt big representatives of $D_j+tH_j$ at the two endpoints.  They can
be chosen to contain small common positive multiples of finitely many
general ample divisors spanning $N^1(P_j/B)$, as well as a fixed ample
part.  To obtain these supports, subtract the desired sufficiently
small ample sum from the big class, represent the remainder effectively,
and mix this representative into a klt one with a small coefficient.
The endpoint representatives determine a rational segment of klt
boundaries.  Allowing small independent changes in the coefficients
of the spanning divisors embeds it in a rational polytope of klt
boundaries with fixed ample part.  The adjoints remain big after
shrinking these changes: on the segment their relative classes are
$N_j+tH_j$, an effective class plus a big class with $t>0$.

BCHM finiteness for this polytope gives finitely many marked ample
models over $B$ \cite[Corollary~1.1.5]{BCHM}.  Every sufficiently late $P_i$ is one of
them.  In fact $P_j\dashrightarrow P_i$ is small, and the pushes of the
chosen divisor classes still span $N^1(P_i/B)$.  For completeness, a
small birational map between $\Q$-factorial models identifies these
numerical spaces: the pullback difference of a numerically trivial
divisor and its transform on a common resolution is exceptional over
the other model and numerically trivial on its contracted curves;
negativity applied with both signs makes that difference zero.  The
same argument in reverse proves the assertion.  At the parameter
$\lambda_i$ the pushed adjoint is nef.  A sufficiently small rational
change in the spanning coefficient directions adds a relatively ample
class on $P_i$, while staying in the polytope, so the resulting adjoint
is ample there.  Sections of the corresponding divisors agree across
the small map, as they are determined by the same inequalities at
prime divisors.  Thus $P_i$, with its marking from $P_j$, is the ample
model of this perturbed boundary.

Distinct negative steps cannot return to the same marked model.
The pullback difference for $K+D$ is effective at each step and is
nonzero at a genuinely negative step.  If it were zero, the two
pullbacks would agree on a common resolution.  A curve dominating a
negative contracted curve would then have negative intersection on
the old side, whereas its image on the other side is contracted to
the same contraction base and has nonnegative intersection there,
a contradiction.  This gives strict discrepancy improvement at some
valuation.  Subsequent discrepancies cannot decrease, so repetition
of the marked pair is impossible.  The finite list of marked ample
models therefore rules out a positive limiting threshold, proving
\eqref{inc:scaling-zero}.

\smallskip\noindent\emph{Preservation of lc centres.}
The lc centres and their incidences have been retained at every step.
Indeed discrepancy monotonicity and crepancy give
\[
 0\leq A_{D_i}(v)\leq A_{C^\natural}(v),
 \qquad A_{C^\natural}(v)=A_{D_i}(v)+v(N_i).
\]
Together with the equality of initial lc places, this shows that the
lc places are exactly the same and have $v(N_i)=0$.  Their centres
are consequently not contained in $\Supp N_i$.  Every negative
contracted curve is contained in $\Supp N_i$, since an effective
$\Q$-Cartier divisor has nonnegative intersection with a curve not
contained in its support.  The contraction is therefore an isomorphism
over a target neighbourhood of the image of each lc centre's generic point.
To see this, a positive-dimensional projective fibre through a point
outside $\Supp N_i$ would contain a contracted curve through that
point.  A zero-dimensional point of a connected fibre cannot be a
separate component of a larger fibre.  Properness and normality of
the target then give the asserted isomorphism neighbourhood.  The
flip agrees over the same neighbourhood.  For an inclusion of lc
centres, test it at the generic point of the smaller centre; this lies
in that common isomorphism locus.  Inclusions are thus preserved in
both directions, and properness preserves their images in $B$.
Only finite compositions will be used below.

\smallskip\noindent\emph{Movable representatives.}
Fix $j$ after the last divisorial contraction in a putative infinite
sequence, or take its final nef model in the finite case.  Fix also
a relatively ample Cartier divisor $A$ on $P_j$.  We claim that for
every rational $\delta>0$ and every finite list of primes on $P_j$
there is an effective divisor avoiding that list with
\begin{equation}\label{inc:movable}
 E_\delta\sim_{\Q,B}N_j+\delta A.
\end{equation}
In the finite case $N_j$ is nef and the sum is ample, so general
representatives on the affine base give the claim.  In the infinite
case choose $i$ first, sufficiently large that
$\delta A-\lambda_iH_j$ is ample on the fixed model $P_j$.
All maps from $P_j$ to $P_i$ are small.  Choose an ample Cartier
divisor $A_i^+$ on $P_i$.  Once $i$ and this divisor have been fixed,
openness of the ample cone permits a positive rational $\alpha$ such
that
\begin{equation}\label{inc:ample-supplement}
 \delta A-\lambda_iH_j-\alpha(A_i^+)_{P_j}
 \quad\hbox{is relatively ample.}
\end{equation}
The divisor $N_i+\lambda_iH_i+\alpha A_i^+$ is also ample, being nef
plus ample.  On the affine base choose effective representatives of
it and of \eqref{inc:ample-supplement}, avoiding the corresponding
finite lists of primes.  Transform the first back to $P_j$ and add
the second.  Smallness identifies the primes and yields
\eqref{inc:movable}.  The choice of $\alpha$ is made \emph{after} $i$;
no bound on the size of $(A_i^+)_{P_j}$ is required.

\smallskip\noindent\emph{Vanishing of the exceptional error.}
We show first that $N_j$ is vertical over $B$.  On the generic fibre
take a common resolution of $P$, $P_j$, and $G_{\rm gen}$.
The effective Cartier pullbacks of $N_j$ are bounded above by those
of $N$ by \eqref{inc:decreasing-pulls}.  Since $N$ is exceptional
over $G_{\rm gen}$, the order of $N_j$ at every prime of
$G_{\rm gen}$ is therefore zero.  This assertion concerns Cartier
pullback data on a common resolution; it does not assert a morphism
$P_j\to G_{\rm gen}$.

Suppose a horizontal component $F$ of $N_j$ has coefficient $n_F>0$.
Apply \eqref{inc:movable} with $F$ in the avoided list.  On the generic
fibre relative linear equivalence becomes ordinary linear equivalence,
so, for a rational function $h_\delta$ and a positive integer
$m_\delta$,
\[
 E_\delta=N_j+\delta A+
       \frac{1}{m_\delta}\operatorname{div}(h_\delta).
\]
Taking its coefficient at $F$ gives, with $a_F=\ord_F(A)$,
\begin{equation}\label{inc:cancel-F}
 \frac{\ord_F(h_\delta)}{m_\delta}=-n_F-\delta a_F.
\end{equation}
Pull this effective divisor to the common resolution and take its
trace on $G_{\rm gen}$.  The zero trace of $N_j$ gives
\begin{equation}\label{inc:Weil-bound}
 \frac{1}{m_\delta}\operatorname{div}_{G_{\rm gen}}(h_\delta)
       +\delta A_{\rm tr}\geq0,
\end{equation}
where $A_{\rm tr}$ is the fixed Weil trace of the Cartier pullback
data of $A$.  Choose a Cartier divisor $H_G\geq A_{\rm tr}$ on the
normal projective variety $G_{\rm gen}$.  Such a divisor exists:
for a sufficiently large multiple of an ample Cartier divisor $L_G$,
the coherent divisorial sheaf $\mathcal O_{G_{\rm gen}}(mL_G-A_{\rm tr})$
has a nonzero section $h$, and $mL_G+\operatorname{div}(h)$ is a
Cartier majorant.  Consequently the divisor obtained from
\eqref{inc:Weil-bound} by replacing $A_{\rm tr}$ with $H_G$ is
effective and $\Q$-Cartier.  Pulling \emph{this} divisor to the
valuation $F$ gives
\[
 \frac{\ord_F(h_\delta)}{m_\delta}
       \geq-\delta\ord_F(H_G).
\]
Together with \eqref{inc:cancel-F}, this implies
\[
 0<n_F\leq\delta\bigl(\ord_F(H_G)-a_F\bigr)
\]
for arbitrarily small positive rational $\delta$, a contradiction.
This proves verticality without assuming $G_{\rm gen}$ is
$\Q$-factorial or pulling back a Weil divisor that is not Cartier.

Next $N_j$ is very exceptional over $B$: it is vertical, and over
every prime divisor $T\subset B$ some prime divisor dominating $T$
is absent from its support.  A component of $N_j$ dominating a base
prime must dominate a boundary prime, because it is the transform of
a toroidal component of $N$.  For a boundary prime $T$, minimize
$\phi$ over a positive vector on its base axis.  A rational minimizer
exists, and its supporting face has $g=0$.  Its toroidal divisorial
valuation has an lc centre dominating $T$ on $P$, hence on $P_j$.
That centre is not contained in $\Supp N_j$.  The pullback of a local
Cartier equation of $T$ vanishes at its generic point.  A prime
component of this Cartier pullback containing that centre dominates
$T$ and is absent from $N_j$.  This reasoning also applies when
the lc centre has higher codimension; the centre itself need not
be a prime divisor.  For a non-boundary prime $T$, surjectivity and
the same Cartier pullback argument supply a dominating prime, and
no component of $N_j$ dominates $T$.

Choose rational $\delta_k\downarrow0$ and representatives in
\eqref{inc:movable} avoiding all components of $N_j$.  For each such
component $F$, their restrictions to $F$ are effective Cartier
divisors after clearing denominators, with proper closed supports.
Every curve in $F$ over $B$ not contained in the countable union of
these supports satisfies
\[
 (N_j+\delta_k A)\cdot C\geq0\quad\hbox{for all }k,
 \qquad N_j\cdot C\geq0.
\]
This is precisely nefness on very general curves of $F/B$.
The morphism $P_j\to B$ is still a projective contraction: its
generic function field is regular over that of the normal base,
so its finite Stein factor is $B$.  Apply Shokurov's very-exceptional
negativity lemma in the form of \cite[Lemma~3.3]{BirkarFlips} to
\[
 D'=-N_j,\qquad (D')^+=0,\qquad (D')^-=N_j.
\]
Its negative part is very exceptional, and $-D'=N_j$ is nef on
the required very general curves.  The lemma gives $D'\geq0$.
Since $N_j\geq0$, it follows that $N_j=0$.

This vanishing occurs on the fixed \emph{finite} model $P_j$.
It makes $K_{P_j}+D_j$ relatively rationally linearly trivial and
rules out any subsequent negative step.  In particular the proposed
infinite sequence cannot occur.  The effective dlt structure and
the asserted centre correspondence are therefore obtained after
finitely many steps, as claimed.
\end{proof}

\subsection{Incidence in a prescribed connected component}

\begin{proposition}[Lc incidence]\label{prop:lc-incidence}
Let $s$ have nonzero base cone $\tau_s$.  Let $\delta_0$ be a source
face on which $g$ vanishes, let $V=\overline{O_{\delta_0}}$ be its
closed stratum, and let $C_s$ be any connected component of $V_s$.
The zero face is allowed, with $V=Y$.  Then there is a source cone
$\beta$ such that
\[
 g|_\beta=0,\qquad p_*(\beta)=\tau_s,\qquad
 \overline{O_\beta}\subset V,\qquad O_\beta\cap C_s\ne\varnothing.
\]
\end{proposition}

\begin{proof}
First disregard $V$ and $C_s$.  A minimizer for $\phi$ over a rational
vector in $\relint(\tau_s)$ has a support face on which $g$ vanishes.
The onto-cone refinements imply that this face maps onto $\tau_s$.
Its stratum occurs over $s$: the smooth stratum map has dense open
image, properness gives a specialization over $s$, and the toroidal
face chart at that specialization has open stratum dense in its
ordinary fibre branch.  This is also the stratum availability used
in \eqref{trait:T}.  Since $Y_s$ is connected, this proves the assertion
when $V=Y$.

Assume $\delta_0\ne0$.  Closed strata in this strict toroidal model
are normal and are unions of strata.  At a deeper stratum the boundary
divisors containing $V$ determine a unique face and a unique normal
orbit-closure germ.  Strictness identifies these divisors globally,
so different local face branches of the same closed stratum cannot
be confused.  In particular their individual ray directions, and
coefficients on those directions, agree at incidences.

We separate the chosen component by a pointed \'{e}tale neighbourhood
\[
 (B_2,s_2)\longrightarrow(B,s),\qquad k(s_2)=k(s)=\kappa.
\]
Indeed the Stein factor $S_V$ of the proper map $V\to B$ is finite
over $B$, also when $V\to B$ is nondominant.  The points of its fibre
over $s$ correspond to the connected components of $V_s$.  Over
the henselization at $s$, the finite algebra splits by the idempotents
separating these points.  The finitely many idempotents descend to
some pointed \'{e}tale neighbourhood.  The associated open-and-closed
piece of $V\times_B B_2$ has precisely $C_s$ in its pointed fibre.
This pullback is normal, and its irreducible components are disjoint.
Thus $C_s$ lies in one such component $V_2$, whose fibre at $s_2$
is exactly $C_s$ as a topological space.  The component $V_2$ is
again a closed stratum.

Shrink to a connected smooth affine neighbourhood of $s_2$.
The pulled-back $Y$ remains integral: its generic fibre is geometrically
integral and its normal \'{e}tale pullback has no additional
vertical component.  It remains strict toroidal with connected fibres.
The forms, cones, and gap function pull back, and the pointed ordinary
fibres are identified.  Use the notation $Y/B,s$ for this new family
during the existence argument.  No lift of the trait to $B_2$ is
needed, and no bound on the degree of this neighbourhood will be used.

Apply Lemma~\ref{inc:effective-model} to obtain the effective dlt
crepant contraction $(P_j,D_j)\to B$.  On its initial toroidal
resolution $P$, one lc centre maps onto $V_2$: take the centre of
a toroidal lc valuation with weights in $\relint(\delta_0)$.
Another lc centre maps onto the closure of an all-zero cone over
$\tau_s$ furnished by the first paragraph.  The latter centre has
base image $\overline{O_{\tau_s}}$.  Transform both to $P_j$.
Inside the second centre choose an lc centre $Z$ minimal by inclusion
among those whose image contains $s$.  This is also inclusion-minimal
among \emph{all} lc centres with image containing $s$: any strictly
smaller one would be a subcentre of the second centre as well.

Koll\'{a}r's linking theorem \cite[Theorem~10]{KollarSources} now
applies.  Its hypotheses hold: the morphism is proper with connected
fibre over $s$, the boundary is effective, the pair is dlt, and
$K_{P_j}+D_j\sim_{\Q,B}0$.  It yields inside the first centre a
subcentre $Z'$ whose image contains $s$ and which is $\mathbb P^1$-linked
to $Z$.  Directly linked centres have the same base image, by the
definition of a link \cite[Definition~9]{KollarSources}; hence so do
centres joined by a chain of links.  Consequently
\[
 s\in f(Z')=f(Z)\subset\overline{O_{\tau_s}}.
\]

Return these centre inclusions through the finite MMP to $P$, using
Lemma~\ref{inc:effective-model}, and then map to $Y$.  On the SNC
model the centres are intersections of coefficient-one divisors,
so are closures of all-zero strata.  A toroidal subdivision maps
such a closure to a stratum closure with $g=0$: a relative interior
zero vector lies in its old support cone, and nonnegativity and
linearity force $g$ to vanish on that cone.  We obtain an all-zero
closed stratum contained in $V_2$, whose image contains $s$ and is
contained in $\overline{O_{\tau_s}}$.  Its image is itself the closure
of a base stratum, by smoothness of the open stratum map and
properness.  Since $s\in O_{\tau_s}$, these two containments force
its image to equal $\overline{O_{\tau_s}}$.

Its cone therefore maps onto $\tau_s$, and its open stratum meets
the ordinary fibre over $s$.  Explicitly, at a closure point in
that fibre the relevant face branch already kills every base
boundary equation; the remaining base parameters are smooth
stratum parameters and can be eliminated.  The open stratum is
dense in this fibre branch, so it supplies the required point.
The point belongs to $(V_2)_s=C_s$.  Projecting through the
pointed \'{e}tale neighbourhood returns an all-zero stratum and
this component incidence to the original ordinary fibre.  This
proves the proposition.
\end{proof}

\subsection{Uniformly short labelled paths}

Return to the original parameter fibre with its bounded Kummer data
and the labelled boundary-weight vector $x$ of $w$.
Write $x$ in the primitive ray coordinates of its simplicial support
cone $\delta$.  Retain the coefficients on rays where $g=0$ and set
the others to zero; denote the resulting vector by $x_0$, its support
face by $\delta_0$, and $g(x)$ by $h$.  Finiteness of the ray data
gives a constant $C_0$, depending only on the bounded families, such that
\begin{equation}\label{inc:projection}
 \begin{gathered}
 \|x-x_0\|\leq C_0h,\qquad b-p_*x_0\geq 0,\\
 \|b-p_*x_0\|\leq C_0h.
 \end{gathered}
\end{equation}
Indeed each discarded coefficient is at most $h/g(r)$ for its primitive
ray $r$, the positive numbers $g(r)$ belong to a fixed finite list, and
the ray maps belong to a fixed finite list as well.  The middle
inequality is coordinatewise in the base cone.  We may use the sum of
absolute ray coefficients for all norms; the finitely many changes of
cone coordinates only change uniform constants.

We construct a path
from the labelled $x$ to a labelled vector defining a valuation $w_*$
with $A_\sigma(w_*)=0$.  If $h=0$, take $w_*=w_x$ and no path.
If the base cone is zero and $h>0$, then $b=0$ and the segment from
$x$ to $x_0$ stays in the same slice and labelled chart.  Its length
is at most $C_0h$, and its endpoint has zero gap.  These cases include
the zero face and all zero-dimensional cells.

It remains to consider $h>0$ and $\tau_s\ne0$.  Put
$V=\overline{O_{\delta_0}}$, taking $V=Y$ if $\delta_0=0$, and
let $C_s$ be the connected component of $V_s$ containing the label
of $x$.  Proposition~\ref{prop:lc-incidence} supplies an all-zero
cone $\beta$ whose labelled stratum meets $C_s$.

Partition $C_s$ by the irreducible components it contains of
$O_\gamma\cap Y_s$, where the stratum closure lies in $V$ and
$p_*(\gamma)=\tau_s$.  These smooth stratum fibres cover $C_s$.
There is a uniform bound $M$ for the number of their irreducible
components.  In fact the strata and their maps belong to the fixed
finite list of finite-type families.  Cover their source and base
by finitely many affine charts and embed the source charts in fixed
affine spaces over the base charts.  The equations have bounded degree,
so B\'{e}zout bounds the number of geometric components of every chart
fibre, and hence of each stratum fibre.  This bound concerns the original
family before the unbounded \'{e}tale neighbourhood used only for
the existence proof.

Make a graph on these components, joining two when one meets the
closure of the other in $C_s$.  The graph is connected.  Otherwise
the unions belonging to two groups of connected graph components
would both be closed: any point of the closure of one piece belongs
to another piece that would then be joined to it.  The finite unions
would partition the connected space $C_s$ into disjoint nonempty
closed subsets.  A simple path therefore joins the starting label
to a label of $\beta$ with at most $M-1$ edges.  At each edge the
strata are comparable by a face inclusion.  A point of the deeper
stratum in the closure of the other component supplies a chart
realizing exactly that face label.

Every cone $\gamma$ on this path contains the face for $V$, with
the same ray directions, so we may transport $x_0$ to it.  Surjectivity
of $\gamma\to\tau_s$ and the fixed ray maps allow us to add weights
on rays above each base axis whose sum projects to $b-p_*x_0$.
Their total size is at most a uniform multiple of $h$, by
\eqref{inc:projection}.  This gives a rational lift of $b$ within
uniform distance of $x_0$.  Because $b$ is interior to $\tau_s$,
there is also a rational lift in $\relint(\gamma)$.  Interpolating
with a sufficiently small positive rational coefficient gives a
vector $x_\gamma\in\relint(\gamma)$ with
\[
 p_*x_\gamma=b,\qquad
 \|x_\gamma-x_0\|\leq C_1h.
\]
The coefficient of interpolation may depend on the chosen lift;
the bound does not.  Use $x$ itself for the starting choice.

For each edge join its two chosen vectors by the segment in the
larger cone, using the chart at the incidence point.  Their common
face label is the specified component, and the labelled valuation
compatibility in Proposition~\ref{prop:trait-comparison} identifies
the evaluations when moving to the next chart.  Each segment has
length at most a fixed multiple of $h$, so the total length is at
most $C_2h$.  The endpoint lies in $\beta$, has gap zero, and has
excess zero by \eqref{trait:TE}.  Its rational monomial valuation
$w_*$ is divisorial, restricts to $u$, and satisfies
$A_\sigma(w_*)=0$.

For $\omega\in\{\psi,\chi\}$ take its comparison ratio $H_\omega$
with $\sigma$ at a common tensor index $p_\omega$, so that
\[
 A_\omega(v)=A_\sigma(v)+\frac{v(H_\omega)}{p_\omega}.
\]
The two estimates \eqref{trait:TV}, first for retraction and then
along the finitely many labelled segments, give in every case
\begin{equation}\label{inc:ratio-path}
 \left|\frac{w(H_\omega)}{p_\omega}
       -\frac{w_*(H_\omega)}{p_\omega}\right|
 \leq C_3\bigl(E(w)+h\bigr)=C_3 A_\sigma(w).
\end{equation}
The constant is uniform because both additional form boundaries are
effective on $G$, which is exactly the hypothesis needed for the
uniform ratio estimates.  Neither the number of auxiliary MMP steps
nor the degree of the \'{e}tale neighbourhood enters it.

Write $q_\omega(v)=v(H_\omega)/p_\omega$.  The assumed inequality
at $w_*$ is $q_\chi(w_*)-c q_\psi(w_*)\geq0$.  Thus
\begin{align*}
 A_\chi(w)-cA_\psi(w)
 &= (1-c)A_\sigma(w)+q_\chi(w)-c q_\psi(w)\\
 &\geq-\bigl(|1-c|+(1+c)C_3\bigr)A_\sigma(w).
\end{align*}
This proves part~(3) of Theorem~\ref{thm:bounded-fibre} for nonzero
base restriction with $\bar A_\sigma(u)=0$, with precisely the
required uniform dependence.

\subsection{Trivial base restriction and the dependency chain}

To treat $w|_K=0$, adjoin one variable $q$ to both fields and pull
$G$ back to $K(q)$.  Extend an absolute form $\omega$ of index
$r_\omega$ by wedging each tensor factor with $d\log q$; denote
the resulting index-$r_\omega$ form by $\omega^{\rm ext}$.
For every divisorial valuation $v$ of $L(q)$,
\begin{equation}\label{inc:extra-variable}
 A_{\omega^{\rm ext}}(v)=A_\omega(v|_L)+J(v),
 \qquad J(v)\geq0,
\end{equation}
where the same $J$ works for all the forms under consideration and
the discrepancy at a trivial valuation is defined to be zero.

Here is the local calculation.  If $v|_L$ is nonzero, extract its
prime $E$ on a smooth model and work over its generic point.  On
the product with $\mathbb P^1_q$, the boundary of the extended form
is the pullback of the form boundary together with the two reduced
sections $q=0,\infty$.  Subtracting $A_\omega(v|_L)$ leaves exactly
the discrepancy of $v$ for the product pair with the reduced
pullback of $E$ and these two sections.  That pair is SNC with
coefficients one, so its discrepancy $J(v)$ is nonnegative and
does not depend on $\omega$.  If $v|_L$ is trivial, work over the
generic point of the model instead; the same calculation uses
just the two sections.  This proves \eqref{inc:extra-variable},
including homogeneous scales.

In particular the extended reference form is absolutely lc.
If $A_{\sigma^{\rm ext}}(v)=0$, both $A_\sigma(v|_L)$ and $J(v)$
vanish.  The hypothesis $A_\chi\geq cA_\psi$ therefore passes to
the extensions on the zero set of the extended reference
discrepancy, including trivial restriction to $L$.  The generic
boundaries are still effective, and the geometric Fano and
singularity assumptions on $G$ persist.  Moreover
\[
 \bar A_{\sigma^{\rm ext}}(\ord_{q=0})=0
 \quad\hbox{on }K(q):
\]
nonnegativity gives one inequality, and the lift $\ord_{q=0}$
trivial on $L$ has zero discrepancy and gives the other.

Given the original nonzero $w$ trivial on $K$, extend it by the
Gauss valuation on $L(q)$ with $q$ of weight one.  This is a rational
divisorial valuation: on the product of an extracting prime for
$w$ and the $q$-line it is the monomial valuation with those two
rational weights.  It restricts to $\ord_{q=0}$ on $K(q)$ and has
$J=0$, since it is monomial for the reduced SNC product pair.
The already proved case applied to this valuation yields the
desired inequality at $w$ by \eqref{inc:extra-variable}.  The total
dimension bound has increased by only one, so its constant is still
uniform in the original numerical data.  For the trivial valuation
itself the inequality is $0\geq0$.

Parts~(1) and~(2), proved in Section~\ref{sec:traits}, and the
preceding argument prove
Theorem~\ref{thm:bounded-fibre} in full.  The reductions of
Section~\ref{sec:tower} now prove Theorem~\ref{thm:lifting-mixing};
Section~\ref{sec:integral-jump} consequently proves
Proposition~\ref{prop:integral-jump}.  The valuative and rounding
arguments then give Theorem~\ref{thm:phase}, and the global reduction
of Section~\ref{sec:global} gives Theorem~\ref{thm:main}.
